\documentclass[UTF8,11pt]{ctexart}
\usepackage{amsmath,amssymb,mathtools,geometry,hyperref}
\geometry{a4paper,margin=24mm}
\newcommand{\LeanRef}[1]{\texttt{\detokenize{#1}}}
\title{经典自旋的关联函数与偶数矩：逐步推导及 Lean 对照}
\author{Classical Spin Krylov — 可追溯证明示范}
\date{}
\begin{document}
\maketitle
\noindent\textbf{状态说明：}本文区分\textbf{已形式化的数学步骤}与\textbf{尚未形式化的物理步骤}。
对应 Lean 文件为 \texttt{ClassicalSpinKrylov.lean}。是否通过检查以 GitHub Actions 最新的 \texttt{lake build} 结果为准。
完整的任意阶 Lanczos 系数通项\textbf{尚未}由 Lean 证明。

\section{物理模型（泊松向量场的部分已形式化）}
设固定长度 $S>0$ 的经典自旋满足
\[
 S_x^2+S_y^2+S_z^2=S^2,\quad
 \{S_i,S_j\}=\epsilon_{ijk}S_k,\quad H=\chi S_z^2.
\]
由 $\dot S_i=\{S_i,H\}$ 得
\[
 \dot S_x=-2\chi S_zS_y,\qquad
 \dot S_y= 2\chi S_zS_x,\qquad \dot S_z=0.
\]
定义 $z=S_z/S\in[-1,1]$，令
$S_x=S\sqrt{1-z^2}\cos\phi$，
$S_y=S\sqrt{1-z^2}\sin\phi$。
因为 $\dot z=0$ 且 $\dot\phi=2\chi Sz$，记 $c=2\chi S$，
\[
 S_x(t)=S\sqrt{1-z^2}\cos(\phi+czt).
\]
其中标准泊松结构常数、哈密顿量导数及三条运动方程已在 \texttt{ClassicalSpinKrylov/Dynamics.lean} 中验证；对应详见 \texttt{docs/dynamics\_zh.tex}。但完整 Poisson 流形构造及候选解满足 ODE 的时间导数检验仍未完成。

\section{归一化关联函数（尚未在 Lean 中形式化）}
均匀球面测度写为 $d\mu=dz\,d\phi/(4\pi)$，初始观测量 $A=S_x$。
首先
\[
 \langle S_x^2\rangle=
 \frac{S^2}{4\pi}\int_{-1}^{1}(1-z^2)\,dz
 \int_0^{2\pi}\cos^2\phi\,d\phi
 =\frac{S^2}{4\pi}\frac43\pi=\frac{S^2}{3}.
\]
另一方面
\[
 \begin{aligned}
 \langle S_x(0)S_x(t)\rangle
 &=\frac{S^2}{4\pi}\int_{-1}^1(1-z^2)\,dz
  \int_0^{2\pi}\cos\phi\cos(\phi+czt)\,d\phi\\
 &=\frac{S^2}{4}\int_{-1}^1(1-z^2)\cos(czt)\,dz,
 \end{aligned}
\]
最后一步由
$\cos\phi\cos(\phi+a)=\tfrac12[\cos a+\cos(2\phi+a)]$，
及第二项沿 $\phi$ 积分为零给出。因而
\[
 C(t)=\frac{\langle S_x(0)S_x(t)\rangle}{\langle S_x^2\rangle}
 =\frac34\int_{-1}^1(1-z^2)\cos(czt)\,dz.
\]
对上式做初等积分，还可得
$C(t)=3(\sin(ct)-ct\cos(ct))/(ct)^3$（$ct\ne0$），
在 $ct=0$ 时由极限定义 $C(0)=1$。
\textbf{注意：}这条三角积分恒等式目前尚未 Lean 证明。

\section{偶数矩：本文重点形式化的数学结论}
定义归一化权重 $\rho(z)=\tfrac34(1-z^2)$ 以及
\[
 \mu_{2n}:=\frac34\,c^{2n}\int_{-1}^1(1-z^2)z^{2n}\,dz,
 \qquad n=0,1,2,\dots.
\]
这是 Lean 中 \LeanRef{evenMoment} 的\textbf{直接定义}，
不是先输入解析答案。

\paragraph{步骤 M1：偶次幂积分。}
对任意非负整数 $n$，
\[
 \begin{aligned}
 I_n&=\int_{-1}^1 z^{2n}\,dz
 =\left[\frac{z^{2n+1}}{2n+1}\right]_{z=-1}^{z=1}\\
 &=\frac{1^{2n+1}-(-1)^{2n+1}}{2n+1}
 =\frac{1-(-1)}{2n+1}
 =\frac{2}{2n+1}.
 \end{aligned}
\]
Lean 对应：\LeanRef{odd_power_neg_one} 和
\LeanRef{even_power_integral}，底层依赖 Mathlib 的
\texttt{intervalIntegral.integral\_pow}。

\paragraph{步骤 M2：分解权重并分别积分。}
由点态恒等式
\[
 (1-z^2)z^{2n}=z^{2n}-z^{2n+2}
\]
及定积分线性性，
\[
 \begin{aligned}
 J_n&:=\int_{-1}^1(1-z^2)z^{2n}\,dz\\
 &=\int_{-1}^1z^{2n}\,dz-\int_{-1}^1z^{2n+2}\,dz\\
 &=\frac2{2n+1}-\frac2{2n+3}\\
 &=\frac{2[(2n+3)-(2n+1)]}{(2n+1)(2n+3)}\\
 &=\boxed{\frac4{(2n+1)(2n+3)}}.
 \end{aligned}
\]
Lean 对应：\LeanRef{spin_weight_integral}；代码逐项调用
积分同余、积分相减与 M1 的结果。

\paragraph{步骤 M3：代入矩的定义。}
由上式逐项得到
\[
 \begin{aligned}
 \mu_{2n}
 &=\frac34 c^{2n}J_n\\
 &=\frac34 c^{2n}\frac4{(2n+1)(2n+3)}\\
 &=\boxed{\frac{3c^{2n}}{(2n+1)(2n+3)}}.
 \end{aligned}
\]
Lean 对应：\LeanRef{evenMoment_formula}。本定理对所有
$n\in\mathbb N$ 和全部实数 $c$ 都有效（包括 $c=0$）。

\paragraph{步骤 M4：归一化和低阶矩。}
\[
 \begin{aligned}
 \mu_0&=\frac3{1\cdot3}=1,\\
 \mu_2&=\frac{3c^2}{3\cdot5}=\frac{c^2}{5},\\
 \mu_4&=\frac{3c^4}{5\cdot7}=\frac{3c^4}{35}.
 \end{aligned}
\]
Lean 对应：\LeanRef{moment_normalized}、
\LeanRef{moment_two}、\LeanRef{moment_four}。

\section{低阶 Lanczos 系数：形式化范围与边界}
对于零对角元的 Lanczos 链，标准矩关系为
$\mu_2=b_1^2$，$\mu_4=b_1^4+b_1^2b_2^2$。
\textbf{该关系尚未从 Lanczos 算符的定义在 Lean 中形式化}；
当前 Lean 仅验证将下述候选式代入后代数恒等式成立。
当 $c\neq0$ 时，
\[
 \begin{aligned}
 b_1^2&=\mu_2=\frac{c^2}{5},\\
 b_2^2&=\frac{\mu_4-b_1^4}{b_1^2}
  =\frac{\frac{3c^4}{35}-\frac{c^4}{25}}{\frac{c^2}{5}}\\
 &=\frac{\frac{15c^4-7c^4}{175}}{\frac{c^2}{5}}
  =\frac{8c^2}{35}.
 \end{aligned}
\]
Lean 对应的\textbf{仅代数}检验是
\LeanRef{fourth_moment_identity} 及
\LeanRef{second_lanczos_identity}。
若 $c=0$，动力学静止，Krylov 链在起点终止，不能对
$b_1^2$ 做除法。

\section{待完成的正式验证}
完整形式化还须补足：自旋轨道满足 ODE 的微分证明、球面测度的投影、
三角积分给出的 $C(t)$、从 Liouville 算符定义到 Lanczos 矩关系、
Gegenbauer 正交性与任意阶 $b_n$。
以上尚未证明的步骤\textbf{不应}被称为已通过 Lean 检查。
\end{document}
